\chapter{Заключение}

В рамках настоящей работы была разработана бесплатформенная инерциальная навигационная система на основе низкоточных микромеханических датчиков для путеизмерительной тележки, обеспечивающая контроль геометрических параметров железнодорожного пути при существенном снижении себестоимости и снаряжённой массы по сравнению с существующими решениями. Реализованы алгоритмы фильтрации сигналов и температурной компенсации смещения нуля, навигационный алгоритм счисления координат и ориентации, программное обеспечение микроконтроллера и персонального компьютера, спроектирована и изготовлена печатная плата коммутации и питания. По состоянию на момент подготовки настоящего отчёта ряд задач выполнен полностью, отдельные задачи находятся в стадии завершения, а часть запланирована к выполнению до защиты выпускной квалификационной работы.

\section*{Итоги проделанной работы}

\begin{enumerate}
	\item Проведён анализ физических принципов работы микроэлектромеханических акселерометров и гироскопов, рассмотрены ёмкостные и пьезоэлектрические схемы преобразования, классифицированы механические, оптические и вибрационные гироскопы (глава~\ref{chap:theory}, разделы~\ref{sec:general_info_nt}, \ref{sec:accelerometer}, \ref{sec:gyroscope}); выполнен обзор существующих путеизмерительных систем отечественных и зарубежных производителей, определено место разрабатываемой системы среди аналогов (глава~\ref{chap:review}, разделы~\ref{sec:cnii_4mdu}--\ref{sec:advantages}). Задача выполнена.
	
	\item Разработано низкоуровневое программное обеспечение микроконтроллера STM32F3 для сбора данных с инерциальных датчиков L3GD20 и LSM303DLHC, их первичной обработки и передачи на персональный компьютер по интерфейсу USB; реализованы драйверы датчиков, циклический опрос сенсоров и протокол двустороннего обмена телеметрией и управляющими командами (раздел~\ref{sec:firmware}). Задача выполнена.
	
	\item Разработаны и реализованы алгоритмы фильтрации сигналов и температурной компенсации смещения нуля. Применён двухступенчатый статистический фильтр на основе метода Уэлфорда с отбраковкой выбросов по правилу $2\sigma$ и кусочно-линейная температурная компенсация по экспериментально определённым калибровочным коэффициентам; алгоритмы обеспечивают подавление высокочастотного шума и систематического температурного дрейфа в условиях ограниченных вычислительных ресурсов микроконтроллера (раздел~\ref{sec:filtering}). Задача выполнена.
	
	\item Разработан навигационный алгоритм счисления координат и ориентации тележки на основе показаний микромеханических датчиков. Обоснован выбор системы координат ENU и аппарата кватернионов, описана процедура начальной выставки по статическому буферу с учётом вращения Земли, сформулирован алгоритм численного интегрирования угловой скорости и ускорения с линейной коррекцией дрейфа (глава~\ref{chap:navigation}). Теоретическое описание алгоритма завершено; программная реализация на персональном компьютере находится в стадии разработки.
	
	\item Создано программное обеспечение для персонального компьютера, включающее графический интерфейс пользователя на базе фреймворка PyQt5 для управления измерительным процессом и модуль постобработки записанных данных на основе библиотек NumPy и Matplotlib (разделы~\ref{sec:gui}, \ref{sec:postprocessing}). Задача выполнена.
	
	\item Спроектирована печатная плата коммутации и питания для подключения внешних модулей к микроконтроллеру STM32F3 DISCOVERY в среде KiCad; плата изготовлена в ООО \enquote{Резонит}, на ней установлены электронные компоненты, проверена работоспособность и соответствие проектным требованиям (раздел~\ref{sec:pcb}). Задача выполнена.
	
	\item Разработка защитного корпуса прибора ведётся в системе автоматизированного проектирования; на момент подготовки отчёта работа находится в стадии конструирования.
	
	\item Натурные испытания разработанной системы и оценка достижимой точности счисления координат и ориентации запланированы к проведению до защиты выпускной квалификационной работы.
\end{enumerate}

\section*{Задачи, планируемые к выполнению до защиты ВКР}

К моменту защиты выпускной квалификационной работы планируется завершить программную реализацию навигационного алгоритма на персональном компьютере, обеспечив сквозную обработку записанных инерциальных данных от первичных показаний датчиков до восстановленной траектории движения. Разработка защитного корпуса прибора предполагает выпуск конструкторской документации, изготовление корпуса и сборку системы в целевом исполнении. Завершающим этапом работы станут натурные испытания собранной системы с проведением серии измерительных проездов и оценкой достижимой точности счисления координат и ориентации в сравнении с заданными требованиями. Полученные по результатам испытаний оценки погрешностей будут использованы для подтверждения практической применимости разработанной системы в составе путеизмерительной тележки и для определения направлений её дальнейшего совершенствования.